\documentclass{../yalumba}

\title{Hybrid Intrinsic--Relational Invariants\\for Reference-Free Relatedness Detection:\\Spectral Ecology v10 and the\\Complete Ten-Version Arc}
\author{Ajax Davis}
\date{April 2026}

\setabstract{%
We present \emph{Spectral Ecology v10}, the culminating iteration of a
ten-version research program that has systematically explored every major
paradigm for comparing genomic module ecosystems without alignment,
variant calling, or reference genomes.  The program began with sparse
adjacency graphs (v1, $-0.74\%$ separation, $-26.87\%$ gap) whose
within-sample spectral summaries were too coarse to distinguish kinship
from noise, progressed through dense co-occurrence graphs (v2,
$+1.55\%$, $-22.53\%$) that first demonstrated positive separation,
attempted cross-sample optimal transport (v3, $+4.93\%$, $-54.67\%$)
whose spectacular separation was undermined by catastrophic pairwise
instability, experimented with coherence fields (v4--v5) and persistence
fields (v6) that revealed the cost of multi-scale averaging, achieved
the closest approach to a passing gap via Sinkhorn-regularized
cooperative stability (v7, $+0.87\%$, $-0.61\%$), introduced
holonomy-based gauge-theoretic invariants with self-baseline
deconfounding (v8, $+1.62\%$, $-6.41\%$), and achieved the strongest
separation in program history via purely intrinsic curvature
distributions (v9, $+3.20\%$, $-7.93\%$).  Spectral Ecology v10
recognizes that v9 and v7 are complementary: v9 maximizes separation
(detecting that related pairs are \emph{more similar on average}) while
v7 minimizes the weakest gap (ensuring that \emph{no unrelated pair
scores higher than the weakest related pair}).  The hybrid combines five
components---cooperative stability from v7, curvature score from v9,
Sinkhorn symmetry from v7, and Bhattacharyya-derived curvature and
distortion penalties from v9---in a weighted composite ($0.45 / 0.20 /
0.15 / 0.10 / 0.10$) that achieves the best balanced result: $+1.50\%$
separation with a weakest gap of $-1.66\%$.  On the CEPH~1463
six-member pedigree, all four parent--child pairs rank in the top~5 of
ten pairwise comparisons, with Bhattacharyya coefficients above 0.94 for
all related pairs.  However, the spouse pair (Father--Mother) ranks
third at 0.4196, interleaved among the parent--child pairs (0.4325 to
0.4030), confirming the persistent coverage confound that has
limited every version of Spectral Ecology.  We provide a comprehensive
retrospective of the ten-version arc, organized into three eras
(within-sample, cross-sample, and intrinsic), present the full hybrid
methodology with weight exploration across three configurations,
analyze component-level correlations to understand which invariants
detect which aspects of inherited structure, and offer an honest
assessment of the program's achievements and limitations.  The gap
trajectory from $-26.87\%$ through $-0.61\%$ (v7) to $-1.66\%$ (v10)
demonstrates that structural invariants from module co-occurrence graphs
contain genuine kinship signal, but that closing the gap to positive
requires a coverage deconfounding strategy that works simultaneously for
both intrinsic and relational invariants---a problem we frame precisely
but do not solve.}

\begin{document}
\maketitle

\tableofcontents
\newpage


% ════════════════════════════════════════════════════════════════════
\section{Introduction}
\label{sec:intro}

\subsection{Two Complementary Failures}

The Spectral Ecology program has spent nine versions searching for
structural invariants of genomic module graphs that distinguish related
from unrelated individuals without alignment or variant calling.  By the
end of version~9, two clear champions had emerged---and they were
champions of \emph{different things}.

\ymarginnote{v9 sees the forest: related pairs are more similar on
average.  v7 tends the weakest tree: no unrelated pair should breach
the related floor.}

Spectral Ecology v7~\cite{cuturi2013sinkhorn}, with Sinkhorn-regularized
optimal transport and 55\% cooperative weighting, achieved a weakest gap
of $-0.61\%$---the closest any spectral method has come to passing.  Its
separation, however, was modest: $+0.87\%$, meaning that the average
related pair scored only marginally higher than the average unrelated
pair.  The signal was narrow but clean.

Spectral Ecology v9, with purely intrinsic curvature distributions
derived from constraint compatibility tensors, achieved $+3.20\%$
separation---the \textbf{best in program history}, nearly matching
Module Persistence v3's $+3.36\%$.  But its weakest gap was $-7.93\%$,
meaning that at least one unrelated pair scored substantially higher
than at least one related pair.  The signal was strong but uneven.

\begin{pullquote}
v9 maximizes the signal; v7 minimizes the noise.\\
v10 asks: can we have both?
\end{pullquote}

The answer, as we show in this paper, is a qualified yes.  By combining
the intrinsic curvature features of v9 with the relational cooperative
stability features of v7 in a five-component weighted score, v10
achieves $+1.50\%$ separation and $-1.66\%$ gap---neither the best
separation nor the best gap, but a combination that no single prior
version achieved.

\subsection{Why Combination Is Non-trivial}

Naively averaging two scoring systems risks inheriting the worst of both.
If v9's high-scoring unrelated pairs are the same pairs as v7's, then
no weighting will help.  If they are different pairs, then a weighted
combination can suppress each method's outliers.

The critical empirical finding is that v9 and v7 fail on \emph{partially
overlapping but distinct} pair sets.  Both struggle with the spouse pair
(Father--Mother), which consistently scores high due to similar
sequencing coverage producing similar module graph structure.  But v9
additionally inflates cross-grandparent pairs (Pat.~GF $\leftrightarrow$
Mat.~GF) via Bhattacharyya overlap, while v7 inflates in-law pairs
(Mat.~GF $\leftrightarrow$ Father) via cooperative stability leakage.
The hybrid can partially suppress both failure modes, though neither
is fully eliminated.

\subsection{Paper Organization}

This paper serves a dual purpose.  First, it presents the Spectral
Ecology v10 hybrid algorithm with complete results on the CEPH~1463
pedigree (\S\ref{sec:methods}--\S\ref{sec:results}).  Second, it
provides a comprehensive retrospective of the entire ten-version arc,
organized into three eras and analyzing the trajectory of ideas that led
from sparse graphs to hybrid invariants
(\S\ref{sec:arc}--\S\ref{sec:discussion}).  The goal is not merely to
report the v10 result---which is respectable but does not pass---but to
extract the lasting lessons from a sustained research program in
structural genomics.


% ════════════════════════════════════════════════════════════════════
\section{The Ten-Version Arc}
\label{sec:arc}

\subsection{Overview}

Table~\ref{tab:arc} presents the complete trajectory of the Spectral
Ecology program across all ten versions, organized by the three eras
that emerged as the program progressed.

\begin{table}[htbp]
\centering
\tablestyle
\caption{The complete Spectral Ecology trajectory on CEPH~1463.
Separation is the percentage by which average related exceeds average
unrelated.  Gap is the percentage by which the weakest related pair
exceeds the strongest unrelated pair (positive = passing).  Bold
indicates best-in-class for that metric.}
\label{tab:arc}
\small
\begin{tabular}{@{}clrrll@{}}
\toprule
\rowcolor{table-header}
\textcolor{white}{\textbf{Ver.}} &
\textcolor{white}{\textbf{Method}} &
\textcolor{white}{\textbf{Sep. (\%)}} &
\textcolor{white}{\textbf{Gap (\%)}} &
\textcolor{white}{\textbf{Era}} &
\textcolor{white}{\textbf{Key Insight}} \\
\midrule
v1 & Sparse graphs & $-0.74$ & $-26.87$ & Within & Adjacency too sparse \\
v2 & Co-occurrence & $+1.55$ & $-22.53$ & Within & First positive separation \\
\midrule
v3 & Optimal transport & $+4.93$ & $-54.67$ & Cross & High sep, catastrophic gap \\
v4 & Coherence fields & $+0.02$ & $-4.99$ & Cross & Spouse pair breaking \\
v5 & Multi-scale & $-0.74$ & $-11.29$ & Cross & Averaging kills signal \\
v6 & Persistence fields & $-0.51$ & $-3.21$ & Cross & Best gap at that time \\
v7 & Sinkhorn + coop & $+0.87$ & $\mathbf{-0.61}$ & Cross & \textbf{Closest to passing} \\
v8 & Holonomy self-base & $+1.62$ & $-6.41$ & Cross & Top-3 all parent--child \\
\midrule
v9 & Intrinsic curvature & $\mathbf{+3.20}$ & $-7.93$ & Intrinsic & \textbf{Best separation} \\
v10 & Hybrid & $+1.50$ & $-1.66$ & Intrinsic & Best balanced \\
\bottomrule
\end{tabular}
\end{table}

\subsection{Era I: Within-Sample Summaries (v1--v2)}

The first era treated each sample in isolation, computing spectral
summaries of module graphs and comparing these scalar or vector
descriptors.

\textbf{v1 (Sparse Adjacency Graphs).}  The original spectral ecology
prototype constructed module co-occurrence graphs using sequential
adjacency on reads---two modules were connected if they appeared on
the same read in sequence.  The resulting graphs were extremely sparse:
most modules had degree one or zero, and the Laplacian spectrum was
dominated by trivial eigenvalues.  The heat kernel signature collapsed
to near-uniform vectors, and ecological role assignment produced
indistinguishable role distributions across all samples.  Result:
$-0.74\%$ separation, $-26.87\%$ gap.  The method was worse than random.

\ymarginnote{v1's failure was not a failure of spectral methods---it was
a failure of graph construction.}

\textbf{v2 (Dense Co-occurrence Graphs).}  The critical insight was that
module co-occurrence should be measured statistically, not sequentially.
By counting co-occurrence of module pairs across all reads and normalizing
by expected frequency under independence, v2 produced dense, weighted
graphs whose structure reflected genuine statistical dependencies.  The
Laplacian spectrum became informative: different samples produced
different eigenvalue density profiles, and heat kernel signatures at
multiple timescales captured multi-resolution structure.  Result:
$+1.55\%$ separation, $-22.53\%$ gap.  The first positive separation
in the program, confirming that spectral structure contains kinship
signal---but the gap was massive, indicating that the signal was
overwhelmed by per-pair noise.

\begin{keyfinding}
The transition from v1 to v2 established the foundational principle:
\emph{co-occurrence probability, not sequential adjacency}, is the
correct edge weight for module graphs.  Every subsequent version builds
on this principle.
\end{keyfinding}

\subsection{Era II: Cross-Sample Mappings (v3--v8)}

The second era introduced cross-sample computation: constructing
explicit mappings between two samples' module graphs and evaluating
the quality of the mapping as a similarity measure.

\textbf{v3 (Optimal Transport).}  Treating each sample's heat kernel
signature distribution as a measure on module space, v3 used optimal
transport (Earth Mover's Distance) to compare samples.  This produced
spectacular separation ($+4.93\%$) but catastrophic gap ($-54.67\%$):
the transport plan was dominated by a few large-module matchings that
happened to align well for certain unrelated pairs.  The lesson was
that raw optimal transport without regularization is too sensitive to
graph structure to serve as a kinship metric.

\textbf{v4 (Coherence Fields).}  Rather than global transport, v4
constructed a point-to-point \emph{coherence field} mapping each module
in sample $A$ to its most compatible counterpart in sample $B$.  The
coherence of this field---measured by smoothness and consistency---served
as the similarity score.  This dramatically reduced the gap to $-4.99\%$
but also reduced separation to $+0.02\%$: the coherence field was too
smooth, assigning high scores to any pair with similar graph density.
The spouse pair broke the ranking entirely.

\textbf{v5 (Multi-Scale Coherence).}  Hypothesizing that single-scale
coherence was too coarse, v5 computed coherence fields at multiple
heat kernel timescales and aggregated them.  The result was worse:
$-0.74\%$ separation, $-11.29\%$ gap.  Averaging across scales
destroyed the narrow kinship signal that existed at individual scales,
confirming that multi-scale analysis requires careful fusion, not
simple averaging.

\textbf{v6 (Persistence Fields).}  Rather than averaging, v6 tracked
which coherence patterns \emph{persisted} across timescales, weighting
them by their persistence.  This was inspired by persistent homology:
features that survive across scales are structural; features that appear
at only one scale are noise.  The gap improved to $-3.21\%$---the best
at that time---but separation was negative ($-0.51\%$), meaning the
method could not even detect that related pairs were more similar on
average.

\ymarginnote{v3--v6: each version improved one metric at the expense of
the other, revealing a fundamental tradeoff.}

\textbf{v7 (Sinkhorn + Cooperative Stability).}  The breakthrough version
introduced two innovations simultaneously.  First, Sinkhorn regularization
of the transport plan~\cite{cuturi2013sinkhorn} stabilized the matching
by enforcing approximate doubly-stochastic structure, preventing the
degenerate matchings that plagued v3.  Second, ``cooperative stability''
measured the degree to which two samples' modules co-supported each other
in a mutual reinforcement framework: modules that were important in both
samples received high cooperative weight.  With 55\% weight on
cooperative stability, v7 achieved $-0.61\%$ gap---the closest any
spectral method has come to passing.  Separation was modest ($+0.87\%$)
because the Sinkhorn regularization compressed the score range, but the
compression was \emph{uniform}, which is why the gap was so tight.

\begin{keyfinding}
v7 demonstrated that regularization is more important than expressiveness:
a constrained transport plan with limited freedom outperformed more
expressive but less stable methods.  The $-0.61\%$ gap remains the
single closest approach to passing in the Spectral Ecology program.
\end{keyfinding}

\textbf{v8 (Holonomy Self-Baseline).}  Drawing on differential geometry
and gauge theory~\cite{nakahara2003geometry}, v8 recast ecosystem
comparison as parallel transport around triangles in a joint module
graph.  The holonomy---the failure of a vector to return to its starting
point after transport around a closed loop---measures the ``curvature''
of the connection between two samples.  A self-baseline (transporting
each sample through itself) provided deconfounding.  The result was
$+1.62\%$ separation with the first correct top-3 ranking (all
parent--child), but the gap regressed to $-6.41\%$ due to NA12889's
anomalously low self-curvature.

\subsection{Era III: Intrinsic Invariants (v9--v10)}

The third era abandoned cross-sample mappings entirely, characterizing
each sample by its internal geometric structure and comparing these
characterizations.

\textbf{v9 (Intrinsic Curvature).}  The paradigm shift: instead of
mapping sample $A$ to sample $B$, v9 computed a curvature distribution
for each sample from its own constraint compatibility tensors (edge
compatibility scores composed around triangles) and compared distributions
directly via Bhattacharyya coefficient and Earth Mover's Distance.
No cross-sample computation at comparison time.  Separation reached
$+3.20\%$---the program best---but gap was $-7.93\%$ because the
Bhattacharyya coefficient is inherently sensitive to graph density
similarity.

\textbf{v10 (Hybrid).}  The present work.  Recognizing that v9 and v7
optimize complementary objectives, v10 combines five components: v7's
cooperative stability and Sinkhorn symmetry with v9's curvature score,
and two Bhattacharyya-derived penalty terms from v9's comparison phase.
The result is $+1.50\%$ separation and $-1.66\%$ gap---neither the best
at either metric, but the best combination of both.

\begin{definition}
\textbf{The separation--gap tradeoff.}  In all Spectral Ecology
versions, improvements in average separation (how far apart the class
means are) tend to come at the expense of the weakest gap (how well
the tails are separated).  This is because methods that increase mean
separation tend to do so by amplifying already-large differences, which
widens the variance within each class.
\end{definition}


% ════════════════════════════════════════════════════════════════════
\section{Methods}
\label{sec:methods}

\subsection{The v9 Component: Intrinsic Curvature}

The intrinsic curvature component operates entirely within each sample.
Given a module co-occurrence graph $G = (V, E, w)$ where vertices are
modules and edge weights $w_{ij}$ represent normalized co-occurrence
frequency, we compute the following quantities.

\subsubsection{Edge Compatibility}

For each edge $(i,j) \in E$, the \emph{compatibility score} $\psi(i,j)$
measures the consistency of the ecological patch tensors at modules $i$
and $j$:
\begin{equation}
\psi(i,j) = \frac{w_{ij}}{\sqrt{w_{ii} \cdot w_{jj}}} \cdot
\left(1 - \frac{|d_i - d_j|}{d_i + d_j}\right)
\label{eq:compatibility}
\end{equation}
where $d_i = \sum_k w_{ik}$ is the weighted degree of module $i$.  The
first factor normalizes edge weight by the geometric mean of self-loops
(equivalent to the cosine similarity of co-occurrence profiles when the
graph is complete), and the second factor penalizes degree asymmetry.

\subsubsection{Triangle Curvature}

For each triangle $(i,j,k)$ in $G$, the \emph{triangle curvature} is:
\begin{equation}
\Delta_{ijk} = \bigl|\psi(i,j) + \psi(j,k) - \psi(i,k)\bigr|
\label{eq:curvature}
\end{equation}

This measures the failure of compatibility to be additive around the
triangle.  In a flat metric space, the triangle inequality is satisfied
with equality for collinear points; deviations from additivity indicate
curvature.  The distribution of $\Delta$ over all triangles in $G$ is the
\emph{curvature distribution} of the sample.

\begin{definition}
\textbf{Curvature distribution.}  For a module co-occurrence graph $G$
with triangle set $\mathcal{T}$, the curvature distribution is the
empirical measure $\mu_G = \frac{1}{|\mathcal{T}|} \sum_{(i,j,k) \in
\mathcal{T}} \delta_{\Delta_{ijk}}$.  This is an intrinsic invariant:
it depends only on the internal structure of $G$, not on any external
reference or mapping.
\end{definition}

\subsubsection{Distribution Comparison}

Two curvature distributions $\mu_A$ and $\mu_B$ are compared via:

\textbf{Bhattacharyya coefficient:}
\begin{equation}
BC(\mu_A, \mu_B) = \sum_{b=1}^{B} \sqrt{h_A(b) \cdot h_B(b)}
\label{eq:bhattacharyya}
\end{equation}
where $h_A$ and $h_B$ are histogram approximations over $B$ bins.

\textbf{Earth Mover's Distance:}
\begin{equation}
\text{EMD}(\mu_A, \mu_B) = \sum_{b=1}^{B} \left|F_A(b) - F_B(b)\right|
\label{eq:emd}
\end{equation}
where $F_A$ and $F_B$ are cumulative distribution functions.

\subsection{The v7 Component: Cooperative Stability}

The cooperative stability component requires cross-sample computation.
Given two samples $A$ and $B$, it proceeds as follows.

\subsubsection{Sinkhorn-Regularized Transport}

We construct a cost matrix $C \in \mathbb{R}^{|V_A| \times |V_B|}$
where $C_{ij}$ is the dissimilarity between the heat kernel signatures
of module $i$ in $A$ and module $j$ in $B$.  The Sinkhorn-regularized
transport plan~\cite{cuturi2013sinkhorn} is:
\begin{equation}
P^* = \operatorname{diag}(u) \, K \, \operatorname{diag}(v)
\label{eq:sinkhorn}
\end{equation}
where $K_{ij} = \exp(-C_{ij}/\varepsilon)$ is the Gibbs kernel with
regularization $\varepsilon > 0$, and $(u,v)$ are the Sinkhorn scaling
vectors obtained by alternating row and column normalization until
convergence.

\subsubsection{Cooperative Stability Score}

The cooperative stability measures the degree to which the transport plan
reveals mutual structural reinforcement:
\begin{equation}
\text{coop}(A,B) = \frac{1}{|V_A|} \sum_{i \in V_A} \max_{j \in V_B}
P^*_{ij} \cdot \mathbb{1}\!\left[\max_{k \in V_A} P^*_{kj} = i\right]
\label{eq:coop}
\end{equation}

This counts the fraction of modules in $A$ that are the preferred partner
of their own preferred partner in $B$---a mutual best-match criterion.
Values range from 0 (no mutual matches) to 1 (perfect mutual matching).

\subsubsection{Sinkhorn Symmetry}

The symmetry of the transport plan is:
\begin{equation}
\text{sym}(A,B) = 1 - \frac{\|P^* - (P^*)^\top\|_F}
{\|P^*\|_F + \|(P^*)^\top\|_F}
\label{eq:symmetry}
\end{equation}

When $|V_A| \neq |V_B|$, we compare the marginals of $P^*$ to the
uniform distribution.  High symmetry indicates that the transport plan
treats both directions equivalently, which is expected for truly similar
structures.

\subsection{The Hybrid Composition}

Spectral Ecology v10 combines five components into a single composite
score:
\begin{equation}
\boxed{
S(A,B) = 0.45 \cdot \text{coop} + 0.20 \cdot \text{curv}
+ 0.15 \cdot \text{sym} + 0.10 \cdot (1 - \kappa) + 0.10 \cdot (1 - d)
}
\label{eq:composite}
\end{equation}
where:
\begin{itemize}
\item $\text{coop} = \text{coop}(A,B)$ is the cooperative stability
  from (\ref{eq:coop})
\item $\text{curv}$ is the curvature similarity score combining
  Bhattacharyya coefficient and EMD
\item $\text{sym} = \text{sym}(A,B)$ is the Sinkhorn symmetry
  from (\ref{eq:symmetry})
\item $\kappa$ is the curvature penalty: the absolute difference
  in mean curvature, $|\bar\Delta_A - \bar\Delta_B|$
\item $d$ is the distortion penalty: the normalized difference in
  graph complexity metrics
\end{itemize}

\begin{keyfinding}
The five-component formula achieves its balance by combining one
\emph{relational} quantity (cooperative stability, weight 0.45), one
\emph{intrinsic} quantity (curvature similarity, weight 0.20), one
\emph{transport quality} indicator (symmetry, weight 0.15), and two
\emph{penalty} terms (curvature and distortion differences, weight
0.10 each).  The relational component dominates because it was the
most effective single predictor in v7.
\end{keyfinding}

\subsection{Weight Exploration}

We explored three weight configurations to understand the
sensitivity of the hybrid to its parameters.

\textbf{Configuration 1: Cooperative-heavy (55/30/15).}
Weights: $0.55 \cdot \text{coop} + 0.30 \cdot \text{curv} + 0.15 \cdot
\text{sym}$.  This produced the highest separation ($+2.06\%$) but the
worst gap ($-6.67\%$), because amplifying cooperative stability widened
the variance among related pairs.

\textbf{Configuration 2: Extra cooperative (65/20/15).}
Weights: $0.65 \cdot \text{coop} + 0.20 \cdot \text{curv} + 0.15 \cdot
\text{sym}$.  Even more cooperative weight ($+1.90\%$ separation,
$-6.57\%$ gap) with no improvement over Configuration~1, confirming
diminishing returns from cooperative stability.

\textbf{Configuration 3: Full five-field (45/20/15/10/10).}
The selected configuration.  Including the penalty terms reduced
separation slightly ($+1.50\%$) but dramatically improved the gap
from $-6.67\%$ to $-1.66\%$, because the curvature and distortion
penalties selectively suppress scores for pairs with different graph
structure---which are disproportionately unrelated.

\begin{limitation}
The weight exploration was manual: three configurations selected by
intuition.  A systematic grid search or Bayesian optimization over the
weight space was not performed.  The selected weights may be locally
but not globally optimal.
\end{limitation}


% ════════════════════════════════════════════════════════════════════
\section{Results}
\label{sec:results}

\subsection{CEPH 1463 Pedigree}

All experiments were conducted on the CEPH~1463 six-member pedigree from
the Illumina Platinum Pedigree Datasets~\cite{ceph1463}, consisting of
a paternal grandfather--grandmother pair (NA12889, NA12890), a maternal
grandfather--grandmother pair (NA12891, NA12892), and their married
offspring (NA12877 = Father, NA12878 = Mother).

\begin{pedigree}
\begin{verbatim}
  Pat.GF (NA12889) --- Pat.GM (NA12890)    Mat.GF (NA12891) --- Mat.GM (NA12892)
           |                                          |
           +-------- Father (NA12877) == Mother (NA12878) --------+
                            ||                    ||
                      (married couple -- not blood related)
\end{verbatim}
\end{pedigree}

There are four parent--child related pairs (Pat.~GF $\to$ Father,
Pat.~GM $\to$ Father, Mat.~GF $\to$ Mother, Mat.~GM $\to$ Mother) and
six unrelated pairs (the spouses plus five in-law or cross-grandparent
pairs).

\subsection{Per-Sample Ecosystem Statistics}

Table~\ref{tab:samples} summarizes the per-sample ecosystem construction
statistics and curvature properties.

\begin{table}[htbp]
\centering
\tablestyle
\caption{Per-sample ecosystem statistics on CEPH~1463.  Modules are
connected components of the k-mer co-occurrence graph.  Triangles are
3-cliques in the module co-occurrence graph.  $\bar\Delta$ is the
mean triangle curvature.  Time is the full preparation phase including
module discovery, graph construction, and curvature computation.}
\label{tab:samples}
\begin{tabular}{@{}llrrrl@{}}
\toprule
\rowcolor{table-header}
\textcolor{white}{\textbf{Sample}} &
\textcolor{white}{\textbf{ID}} &
\textcolor{white}{\textbf{Modules}} &
\textcolor{white}{\textbf{Triangles}} &
\textcolor{white}{\textbf{$\bar\Delta$}} &
\textcolor{white}{\textbf{Time (s)}} \\
\midrule
Pat.~GF & NA12889 & 457 & 4{,}822 & 1.064 & 60.8 \\
Pat.~GM & NA12890 & 607 & 1{,}875 & 1.434 & 51.1 \\
Mat.~GF & NA12891 & 712 & 4{,}867 & 1.361 & 59.5 \\
Mat.~GM & NA12892 & 593 & 3{,}714 & 1.051 & 54.3 \\
Father  & NA12877 & 492 & 1{,}256 & 1.108 & 51.7 \\
Mother  & NA12878 & 671 & 3{,}296 & 1.307 & 77.0 \\
\bottomrule
\end{tabular}
\end{table}

Several patterns are noteworthy.  Module counts range from 457 (Pat.~GF)
to 712 (Mat.~GF), reflecting coverage-dependent graph construction.
Triangle counts are \emph{not} monotonic in module count: Pat.~GM has 607
modules but only 1{,}875 triangles, while Pat.~GF has 457 modules but
4{,}822 triangles, indicating very different graph densities.  Mean
curvature $\bar\Delta$ varies from 1.051 (Mat.~GM) to 1.434 (Pat.~GM),
a 36\% relative range.  Preparation time varies from 51.1\,s (Pat.~GM)
to 77.0\,s (Mother), dominated by triangle enumeration.

\subsection{Full Pairwise Comparison}

Table~\ref{tab:pairs} presents all ten pairwise comparisons with the
full five-component decomposition, ranked by composite score.  Related
pairs are marked with $\blacktriangleright$.

\begin{table}[htbp]
\centering
\tablestyle
\caption{All 10 pairwise comparisons on CEPH~1463, ranked by composite
score.  Components: coop = cooperative stability, curv = curvature
score, sym = Sinkhorn symmetry, BC = Bhattacharyya coefficient.
Related pairs marked with $\blacktriangleright$.}
\label{tab:pairs}
\small
\begin{tabular}{@{}clccccr@{}}
\toprule
\rowcolor{table-header}
\textcolor{white}{\textbf{Rank}} &
\textcolor{white}{\textbf{Pair}} &
\textcolor{white}{\textbf{Coop}} &
\textcolor{white}{\textbf{Curv}} &
\textcolor{white}{\textbf{Sym}} &
\textcolor{white}{\textbf{BC}} &
\textcolor{white}{\textbf{Score}} \\
\midrule
$\blacktriangleright$ 1
  & Mat.~GF $\to$ Mother
  & 0.1212 & 0.9053 & 0.8972 & 0.9861 & 0.4325 \\
$\blacktriangleright$ 2
  & Pat.~GM $\to$ Father
  & 0.0971 & 0.8639 & 0.8660 & 0.9549 & 0.4218 \\
3 & Spouses
  & 0.1475 & 0.8987 & 0.8518 & 0.9725 & 0.4196 \\
4 & Mat.~GF $\leftrightarrow$ Father
  & 0.1072 & 0.8521 & 0.8732 & 0.9760 & 0.4158 \\
$\blacktriangleright$ 5
  & Mat.~GM $\to$ Mother
  & 0.1518 & 0.8675 & 0.9034 & 0.9686 & 0.4157 \\
6 & Pat.~GF $\leftrightarrow$ Mat.~GF
  & 0.0901 & 0.8310 & 0.8707 & 0.9662 & 0.4047 \\
$\blacktriangleright$ 7
  & Pat.~GF $\to$ Father
  & 0.0774 & 0.8194 & 0.8228 & 0.9467 & 0.4030 \\
8 & Pat.~GM $\leftrightarrow$ Mat.~GM
  & 0.0814 & 0.7914 & 0.8680 & 0.9342 & 0.4004 \\
9 & Pat.~GF $\leftrightarrow$ Mother
  & 0.0788 & 0.8124 & 0.8373 & 0.9660 & 0.3912 \\
10 & Pat.~GF $\leftrightarrow$ Mat.~GM
  & 0.0675 & 0.8063 & 0.8576 & 0.9725 & 0.3878 \\
\bottomrule
\end{tabular}
\end{table}

\subsection{Ranking Analysis}

\begin{keyfinding}
All four parent--child pairs rank in the top~7 of ten pairs, with three
of the four in the top~5.  The Mat.~GF $\to$ Mother pair ranks first
with a clear margin ($+0.0107$ over rank~2).  The Pat.~GF $\to$ Father
pair ranks seventh---the weakest related pair---dragged down by Pat.~GF's
low cooperative stability (0.0774) and low symmetry (0.8228).
\end{keyfinding}

The ranking reveals a critical structural pattern.  The two strongest
related pairs both involve the Mother (NA12878): Mat.~GF $\to$ Mother at
rank~1 and Mat.~GM $\to$ Mother at rank~5.  Mother's ecosystem has 671
modules and 3{,}296 triangles, giving it rich structure for both
cooperative matching and curvature comparison.  By contrast, the two
weaker related pairs both involve the Father (NA12877), who has only
492 modules and 1{,}256 triangles---the sparsest ecosystem in the
pedigree.  Sparse ecosystems produce lower cooperative stability scores
(fewer mutual best-matches are possible) and noisier curvature
distributions (fewer triangles to average over).

\subsection{The Spouse Confound}

The spouse pair (Father--Mother) ranks third with score 0.4196, nestled
between two parent--child pairs.  This is the \emph{persistent spouse
confound} that has affected every version of Spectral Ecology.

The component decomposition reveals why.  The spouse pair has the
\emph{highest} cooperative stability among all pairs (0.1475), exceeding
every related pair.  This occurs because Father and Mother, as the two
offspring in the pedigree, were sequenced in the same batch with similar
coverage, producing module graphs of compatible size and density.
Compatible graph structure generates compatible heat kernel signatures,
which produce well-matched Sinkhorn transport plans, which inflate
cooperative stability.

\begin{limitation}
The spouse pair's cooperative stability of 0.1475 exceeds the highest
related cooperative stability (Mat.~GM $\to$ Mother, 0.1518) by only
$-2.8\%$, but exceeds the average related cooperative stability (0.1119)
by $+31.8\%$.  Cooperative stability is fundamentally confounded by graph
size compatibility, which correlates with sequencing batch effects.
\end{limitation}

The curvature score partially compensates: the spouse pair's curvature
score (0.8987) is lower than Mat.~GF $\to$ Mother's (0.9053), and its
symmetry (0.8518) is the second-lowest among the top~5.  But the
compensation is insufficient: the cooperative stability advantage is too
large for the curvature and symmetry penalties to overcome.

\subsection{Component Correlation Analysis}

To understand which components carry independent signal, we compute
Pearson correlations between component values across all ten pairs.

\begin{table}[htbp]
\centering
\tablestyle
\caption{Component relationships across the ten pairs.  ``Related
advantage'' shows the percentage by which the average related value
exceeds the average unrelated value.  Higher is better for all
components except curvature penalty ($\kappa$) and distortion ($d$),
where lower is better.}
\label{tab:components}
\begin{tabular}{@{}lcccc@{}}
\toprule
\rowcolor{table-header}
\textcolor{white}{\textbf{Component}} &
\textcolor{white}{\textbf{Avg Related}} &
\textcolor{white}{\textbf{Avg Unrelated}} &
\textcolor{white}{\textbf{Advantage}} &
\textcolor{white}{\textbf{Discriminative?}} \\
\midrule
Cooperative stability & 0.1119 & 0.0954 & $+17.3\%$ & Moderate \\
Curvature score       & 0.8640 & 0.8320 & $+3.8\%$  & Moderate \\
Symmetry              & 0.8724 & 0.8594 & $+1.5\%$  & Weak \\
Bhattacharyya coeff.  & 0.9641 & 0.9646 & $-0.05\%$ & None \\
\bottomrule
\end{tabular}
\end{table}

The analysis reveals a hierarchy of discriminative power.  Cooperative
stability has the largest related advantage ($+17.3\%$) but is also the
most confounded (the spouse pair has the highest value).  Curvature score
has moderate advantage ($+3.8\%$) with less confounding.  Symmetry adds
weak signal ($+1.5\%$).  The Bhattacharyya coefficient is essentially
non-discriminative ($-0.05\%$ advantage), suggesting that histogram
overlap between curvature distributions is too coarse to capture kinship
differences.

\begin{keyfinding}
The Bhattacharyya coefficient, which was the most promising component in
v9 (achieving $BC > 0.97$ for all related pairs), loses its
discriminative power when incorporated into the hybrid.  This is because
the unrelated pairs \emph{also} have high Bhattacharyya coefficients
(average 0.9646 versus 0.9641 for related).  Curvature distribution
overlap is a necessary but not sufficient condition for relatedness.
\end{keyfinding}

\subsection{Weight Configuration Comparison}

Table~\ref{tab:weights} compares the three weight configurations
explored.

\begin{table}[htbp]
\centering
\tablestyle
\caption{Weight configuration comparison.  Sep.\ is average separation.
Gap is the weakest-pair gap.  The selected configuration (3) achieves the
best balanced result.}
\label{tab:weights}
\begin{tabular}{@{}lllrr@{}}
\toprule
\rowcolor{table-header}
\textcolor{white}{\textbf{Config}} &
\textcolor{white}{\textbf{Weights}} &
\textcolor{white}{\textbf{Components}} &
\textcolor{white}{\textbf{Sep.\ (\%)}} &
\textcolor{white}{\textbf{Gap (\%)}} \\
\midrule
1 & 55/30/15       & coop/curv/sym   & $+2.06$ & $-6.67$ \\
2 & 65/20/15       & coop/curv/sym   & $+1.90$ & $-6.57$ \\
\rowcolor{yalumba-cream}
3 & 45/20/15/10/10 & full five-field & $+1.50$ & $-1.66$ \\
\bottomrule
\end{tabular}
\end{table}

Configuration~3's advantage is clear: the two penalty terms reduce
separation by only 0.56 percentage points (from $+2.06\%$ to $+1.50\%$)
but improve the gap by 5.01 percentage points (from $-6.67\%$ to
$-1.66\%$).  This is because the penalty terms selectively suppress
scores for structurally dissimilar pairs, and structurally dissimilar
pairs are disproportionately unrelated.  The curvature penalty $1-\kappa$
punishes pairs with different mean triangle curvature, while the
distortion penalty $1-d$ punishes pairs with different graph complexity.
Both penalties are small for related pairs (who inherit similar graph
structure from shared lineage) and large for unrelated pairs (whose graph
structure depends on independent sequencing conditions).

\subsection{Full Trajectory Comparison}

Table~\ref{tab:trajectory} presents the complete v1--v10 trajectory
alongside comparison methods.

\begin{table}[htbp]
\centering
\tablestyle
\caption{Complete algorithm comparison.  The Spectral Ecology trajectory
is shown alongside the two passing symbiogenesis algorithms and the
gold-standard Rare-Run P90.  Bold indicates best spectral result for
each metric.}
\label{tab:trajectory}
\small
\begin{tabular}{@{}lrrl@{}}
\toprule
\rowcolor{table-header}
\textcolor{white}{\textbf{Algorithm}} &
\textcolor{white}{\textbf{Sep.\ (\%)}} &
\textcolor{white}{\textbf{Gap (\%)}} &
\textcolor{white}{\textbf{Status}} \\
\midrule
Spectral Ecology v1   & $-0.74$  & $-26.87$ & Fail \\
Spectral Ecology v2   & $+1.55$  & $-22.53$ & Fail \\
Spectral Ecology v3   & $+4.93$  & $-54.67$ & Fail \\
Spectral Ecology v4   & $+0.02$  & $-4.99$  & Fail \\
Spectral Ecology v5   & $-0.74$  & $-11.29$ & Fail \\
Spectral Ecology v6   & $-0.51$  & $-3.21$  & Fail \\
Spectral Ecology v7   & $+0.87$  & $\mathbf{-0.61}$  & Best gap \\
Spectral Ecology v8   & $+1.62$  & $-6.41$  & Fail \\
Spectral Ecology v9   & $\mathbf{+3.20}$  & $-7.93$  & Best sep \\
\rowcolor{yalumba-cream}
Spectral Ecology v10  & $+1.50$  & $-1.66$  & Best balanced \\
\midrule
Module Persistence v3  & $+3.36$  & $+0.13$  & \textbf{Pass} \\
Coalition Transfer v3  & $+1.97$  & $+0.06$  & \textbf{Pass} \\
Rare-Run P90           & $+583$   & $+300$   & \textbf{Gold standard} \\
\bottomrule
\end{tabular}
\end{table}

\subsection{Gap Trajectory}

The gap trajectory across all versions is:
\[
-26.87\% \xrightarrow{\text{v2}} -22.53\% \xrightarrow{\text{v4}} -4.99\%
\xrightarrow{\text{v6}} -3.21\% \xrightarrow{\text{v7}} -0.61\%
\xrightarrow{\text{v10}} -1.66\%
\]

\ymarginnote{The gap trajectory is monotonically improving through v7.
v10 is a regression from v7's gap, though still the second-best.}

Note that v3, v5, v8, and v9 are excluded from this sequence because
their gaps were worse than their predecessors---each was an excursion
into higher separation at the cost of gap stability.  The ``gap-improving
subsequence'' (v1, v2, v4, v6, v7) shows steady progress from $-26.87\%$
to $-0.61\%$, with v10 representing a slight regression.

\begin{limitation}
v10's gap of $-1.66\%$ is worse than v7's $-0.61\%$.  The hybrid
inherits curvature-based scores that inflate certain unrelated pairs
beyond what v7's pure cooperative+symmetry scoring produced.  Combining
components introduces new sources of score variance that partially
offset the benefits of multi-signal fusion.
\end{limitation}

\subsection{Timing}

Total pipeline time: 366.3\,s (6 minutes 6 seconds).

\begin{itemize}
\item \textbf{Preparation:} 354.3\,s (96.7\%)---module discovery, graph
  construction, eigendecomposition, heat kernel computation, curvature
  distribution extraction.  This is a per-sample cost, embarrassingly
  parallel.
\item \textbf{Comparison:} 12.1\,s (3.3\%)---10 pairwise comparisons at
  1.21\,s each.  This includes Sinkhorn transport, cooperative stability,
  curvature distribution comparison, and symmetry computation.
\item \textbf{Per-pair marginal cost:} Adding a seventh sample would cost
  ${\sim}60$\,s preparation plus ${\sim}7$\,s comparison (6 new pairs at
  1.21\,s), for ${\sim}67$\,s marginal.
\end{itemize}

The $O(n)$ preparation and $O(n^2)$ comparison structure means that
comparison time dominates only for $n > 250$ samples.  For the
foreseeable cohort sizes ($n < 50$), preparation is the bottleneck.


% ════════════════════════════════════════════════════════════════════
\section{Discussion}
\label{sec:discussion}

\subsection{Why the Hybrid Works Better Than Either Alone}

The fundamental insight of v10 is that v9 and v7 detect \emph{different
aspects} of inherited structure, and their failure modes are partially
non-overlapping.

\textbf{Cooperative stability} (v7) detects relational consistency: do
the modules in sample $A$ have stable, reciprocal partners in sample $B$?
This is a cross-sample quantity that measures how well two ecosystems
can be aligned.  Its failure mode is graph size compatibility: samples
with similar numbers of modules and similar degree distributions produce
well-aligned transport plans regardless of kinship.

\textbf{Curvature similarity} (v9) detects structural isomorphism: do
samples $A$ and $B$ have the same distribution of internal geometric
constraints?  This is an intrinsic quantity that measures whether two
ecosystems have the same ``shape'' without reference to each other.  Its
failure mode is graph density sensitivity: samples with similar edge
density produce similar curvature distributions regardless of kinship.

By combining both, v10 requires that a pair score high on \emph{both}
relational consistency and structural isomorphism.  An unrelated pair
might have compatible graph sizes (high cooperative stability) but
different curvature distributions, or similar curvature distributions
but incompatible transport plans.  The hybrid penalizes both cases.

\begin{pullquote}
v10's principle: related pairs should be similar in structure
\textbf{and} well-matched in organization.  Requiring both is harder
to satisfy by chance than requiring either alone.
\end{pullquote}

\subsection{The Separation--Gap Tradeoff}

The ten-version trajectory reveals a persistent tension between
separation and gap that is not merely an engineering problem but
reflects a structural property of the scoring landscape.

Methods that increase separation do so by amplifying the scores of
already-high-scoring related pairs---typically those involving Mother
(NA12878), whose rich ecosystem (671 modules, 3{,}296 triangles)
provides more evidence for both cooperative stability and curvature
similarity.  But amplifying high scores widens the \emph{variance}
within the related class, because the weakest related pair
(Pat.~GF $\to$ Father) does not benefit proportionally from its
sparse 492-module, 1{,}256-triangle ecosystem.

Methods that minimize the gap do so by compressing the overall score
range---through regularization (Sinkhorn in v7) or through penalty
terms (curvature and distortion in v10).  Compression reduces variance
in both classes, bringing the weakest related score closer to the
strongest unrelated score, but also bringing the class means closer
together.

\begin{definition}
\textbf{Separation--gap duality.}  For a scoring function $S$ and
classes $\mathcal{R}$ (related) and $\mathcal{U}$ (unrelated):
\begin{align}
\text{Sep} &= \frac{\mathbb{E}[S|\mathcal{R}] -
  \mathbb{E}[S|\mathcal{U}]}{\mathbb{E}[S|\mathcal{U}]} \\
\text{Gap} &= \frac{\min_{r \in \mathcal{R}} S(r) -
  \max_{u \in \mathcal{U}} S(u)}{\max_{u \in \mathcal{U}} S(u)}
\end{align}
Separation depends on class means; gap depends on class extrema.
Reducing variance (compression) improves gap at the expense of
separation.  Increasing signal (amplification) improves separation at
the expense of gap.
\end{definition}

v10 navigates this tradeoff by using the amplification of curvature
scores (v9's contribution) with the compression of penalty terms (which
bring outlier unrelated scores down) and cooperative stability (v7's
already-compressed contribution).  The result is a Pareto-optimal
point: no single-method version achieves both $+1.50\%$ separation
\emph{and} $-1.66\%$ gap.

\subsection{What Each Component Detects}

The five components of the v10 score capture distinct aspects of
ecosystem similarity:

\textbf{Cooperative stability} ($w = 0.45$): Measures relational
coherence---whether modules in two samples form reciprocal best-match
partnerships through the Sinkhorn transport plan.  Sensitive to graph
size compatibility and sequencing coverage.  Best at detecting pairs
that share ecological ``roles'' in compatible proportions.

\textbf{Curvature score} ($w = 0.20$): Measures intrinsic structural
similarity---whether two samples' triangle curvature distributions have
similar shape.  Combines Bhattacharyya coefficient (histogram overlap)
and Earth Mover's Distance (distributional displacement).  Sensitive to
graph density and triangle count.  Best at detecting pairs with similar
internal organizational constraints.

\textbf{Sinkhorn symmetry} ($w = 0.15$): Measures transport quality---
whether the optimal alignment between two ecosystems treats both
directions equivalently.  High symmetry indicates that neither sample
has ``orphaned'' modules without partners.  Less confounded by coverage
than cooperative stability because asymmetry is relative, not absolute.

\textbf{Curvature penalty} ($w = 0.10$): Penalizes differences in mean
curvature $\bar\Delta$.  Acts as a filter: pairs with very different
curvature statistics are likely unrelated regardless of other scores.
This is the most discriminative penalty term because mean curvature
varies 36\% across samples.

\textbf{Distortion penalty} ($w = 0.10$): Penalizes differences in
graph complexity metrics.  Complements the curvature penalty by capturing
first-order graph properties (degree distribution, density) that
curvature misses.

\subsection{The Persistent Spouse Confound}

The spouse pair (Father--Mother) ranks third in v10 with score 0.4196,
continuing a pattern that has persisted across all ten versions.  The
spouse confound is not a bug in the algorithm; it is a structural
feature of the dataset.

Father and Mother are the only two offspring in the pedigree.  In the
Illumina Platinum Pedigree sequencing protocol, they were sequenced
under the same conditions, at similar coverage, using the same
library preparation.  This means:

\begin{enumerate}
\item Their module counts are compatible (492 vs.\ 671---within
  the range of parent--child pairs).
\item Their graph densities are compatible (both reflect the same
  library preparation biases).
\item Their curvature distributions are similar (similar sequencing
  artifacts produce similar edge weight patterns).
\item Their cooperative stability is high (compatible graph sizes
  produce well-matched Sinkhorn transport).
\end{enumerate}

\begin{limitation}
The spouse confound is a coverage confound.  Two unrelated individuals
sequenced at similar coverage with similar library preparation will
produce similar module graph statistics across \emph{all} of v10's
components, because all components ultimately derive from the module
co-occurrence graph, which is constructed from k-mer co-occurrence
counts that scale with coverage.  No combination of weights can
eliminate this confound without either (a) explicit coverage
normalization or (b) a feature that is provably coverage-invariant.
\end{limitation}

\subsection{The 0.61\% versus 1.66\% Gap}

A natural question is whether v10 represents progress or regression
relative to v7, given that v7's gap ($-0.61\%$) is closer to passing
than v10's ($-1.66\%$).

The answer depends on what we value.  If the goal is purely to pass
(achieve positive gap), then v7 remains the strongest candidate:
$-0.61\%$ is 2.7$\times$ closer to zero than $-1.66\%$.  If the
goal is to understand the structure of the scoring landscape, then v10
is more informative: it demonstrates that combining intrinsic and
relational invariants improves the separation--gap tradeoff surface,
even if the optimal point on that surface is not the closest approach
to the gap axis.

More importantly, v10 demonstrates that the gap is not sensitive
to the choice of individual invariant but to the \emph{coverage
confound}.  Both v7 and v10 are limited by the same pair
(Pat.~GF $\to$ Father = 0.4030, below the spouse pair at 0.4196),
and both are limited for the same reason (Pat.~GF's sparse ecosystem
produces low cooperative stability).  The gap will not improve
substantially until the coverage confound is addressed directly.

\subsection{Connection to the Theoretical Framework}

The ten-version arc can be understood through the abstraction hierarchy
proposed in the yalumba theoretical framework:

\begin{enumerate}
\item \textbf{Token level:} Shared k-mers, Jaccard similarity.  This is
  what Rare-Run P90 operates on, and it works.
\item \textbf{Summary level:} Per-sample statistics---module counts,
  eigenvalue distributions, degree sequences.  v1--v2 operated here.
\item \textbf{Matching level:} Cross-sample alignments---optimal
  transport, coherence fields, Sinkhorn plans.  v3--v8 operated here.
\item \textbf{Field level:} Spatially organized quantities over module
  graphs---curvature fields, cooperative fields, symmetry fields.  The
  penalty terms in v10 operate here.
\item \textbf{Intrinsic level:} Per-sample geometric invariants---
  curvature distributions, topological features, spectral invariants.
  v9 and the curvature component of v10 operate here.
\end{enumerate}

\ymarginnote{Each level up the hierarchy gains invariance at the cost
of discrimination.  The question is whether kinship signal survives at
the intrinsic level.}

The program's trajectory has been a systematic ascent through this
hierarchy.  Each level offers more invariance (less sensitivity to
nuisance variation) but also less discrimination (less ability to
distinguish genuine signal from noise).  The ten-version result
suggests that kinship signal \emph{does} survive at the intrinsic level
($+3.20\%$ separation in v9), but that the signal-to-noise ratio at
that level is insufficient to separate all pairs ($-7.93\%$ gap).
The hybrid v10 attempts to combine levels 3 and 5, and achieves a
better tradeoff than either alone.

\subsection{Is Spectral Ecology Fundamentally Limited?}

This is the central question the ten-version program raises.  Two
interpretations are possible.

\textbf{Engineering interpretation:} The spectral ecology approach
is correct in principle, but the current implementation is limited by
(a) insufficient coverage deconfounding, (b) too-coarse graph construction
(binary module membership instead of soft assignment), (c) too few
samples for statistical calibration, and (d) too few reads per sample
(50{,}000) for the spectral signal to dominate noise.  Under this
interpretation, a better-engineered version---with explicit coverage
normalization, soft module assignment, larger cohorts, and deeper
sequencing---would eventually pass.

\textbf{Fundamental interpretation:} Module co-occurrence graphs at
the resolution achievable from 50{,}000 reads per sample do not contain
enough structural information to separate first-degree relatives from
coverage-matched unrelated individuals.  The spectral invariants capture
genuine structure, but that structure is dominated by sequencing depth
rather than by inherited genomic organization.  Under this interpretation,
the approach requires a fundamentally different graph construction---
perhaps from raw k-mer co-occurrence without the module abstraction
layer---to access the signal.

\ymarginnote{After ten versions, we believe the truth lies between
these interpretations---closer to the engineering side.}

We note that Module Persistence v3 ($+3.36\%$ sep, $+0.13\%$ gap)
and Coalition Transfer v3 ($+1.97\%$ sep, $+0.06\%$ gap) both pass
using the same module abstraction layer and the same 50{,}000-read
budget.  These methods operate at the token/summary levels of the
hierarchy---they compare which modules are shared and how they
co-occur, not the geometric properties of the co-occurrence graph.
This suggests that the kinship signal is present in the module
co-occurrence data, but that the spectral/geometric invariants
extracted by Spectral Ecology are not the most efficient way to
access it.

\subsection{Limitations}

\begin{enumerate}
\item \textbf{Single pedigree.}  All results are on the CEPH~1463
  six-member pedigree.  Generalization to other pedigrees, populations,
  and sequencing platforms is untested.

\item \textbf{First-degree only.}  The pedigree contains only
  parent--child (first-degree) relationships.  The method's ability to
  detect second-degree (grandparent--grandchild, half-siblings) or
  more distant relationships is unknown.

\item \textbf{Small sample size.}  Six samples produce only ten
  pairwise comparisons (four related, six unrelated).  Statistical
  significance cannot be established from this sample size.

\item \textbf{Manual weight selection.}  The five-component weights
  were chosen from three manually selected configurations.  Systematic
  optimization was not performed.

\item \textbf{Coverage confound unresolved.}  The persistent spouse
  confound demonstrates that all five components are sensitive to
  sequencing coverage.  No coverage-invariant feature was developed.

\item \textbf{50{,}000-read cap.}  The maximum reads per sample was
  50{,}000.  With deeper sequencing, the module co-occurrence graph
  would be denser and the spectral signal might be stronger, but
  this was not tested.

\item \textbf{No cross-validation.}  Weights were selected on the
  same dataset used for evaluation.  This is not a valid evaluation
  protocol for a production system.

\item \textbf{Computational cost.}  At 366\,s for six samples, the
  method is 100$\times$ slower than Rare-Run P90 and produces
  strictly worse results.  Spectral methods are justified only if they
  reveal structure that token methods miss.

\item \textbf{No formal null model.}  The gap and separation metrics
  are computed against an empirical unrelated distribution (six pairs),
  not against a formal null model of unrelated similarity.

\item \textbf{Component correlation.}  The five components are not
  independent: cooperative stability and curvature score are both
  derived from the same module co-occurrence graph.  The effective
  degrees of freedom in the composite score are less than five.
\end{enumerate}


% ════════════════════════════════════════════════════════════════════
\section{Future Work}
\label{sec:future}

The ten-version Spectral Ecology program points toward several
promising directions for future research.

\subsection{Coverage-Invariant Curvature}

The most pressing need is a curvature computation that is provably
invariant to sequencing coverage.  Current curvature distributions
depend on edge weights that scale with coverage, producing similar
distributions for coverage-matched unrelated pairs.  Possible approaches
include:

\begin{itemize}
\item \textbf{Rank-based edge weights:} Replacing co-occurrence
  counts with co-occurrence \emph{ranks} (within each sample) would
  eliminate dependence on absolute coverage while preserving relative
  structure.
\item \textbf{Coverage regression:} Explicitly modeling the
  relationship between coverage and curvature statistics, then
  regressing it out before comparison.
\item \textbf{Spectral normalization:} Normalizing the Laplacian
  spectrum to unit trace before computing heat kernels, making the
  spectral decomposition scale-invariant.
\end{itemize}

\subsection{Gromov--Wasserstein Distance}

Rather than comparing curvature distributions (which are summaries),
compare the full metric structure of module co-occurrence graphs using
Gromov--Wasserstein distance~\cite{memoli2011gromov}.  This is a
distance between metric measure spaces that does not require a common
embedding or point-to-point correspondence, making it a natural fit
for comparing ecosystems of different sizes.

\subsection{Persistent Homology}

The module co-occurrence graph has natural filtration parameters
(edge weight thresholds, heat kernel timescales) that can be used to
construct a persistence diagram~\cite{edelsbrunner2010computational,
kerber2017geometry}.  Persistent homology captures topological
features---connected components, loops, voids---that are invariant
to graph perturbation and potentially more robust to coverage variation
than spectral features.

\subsection{GPU Acceleration}

The preparation phase (354\,s) is dominated by eigendecomposition and
heat kernel computation, both of which are amenable to GPU acceleration
via the yalumba SPIR-V pipeline~\cite{chung1997spectral}.  A
GPU-accelerated version could reduce preparation time to under 10\,s,
making the spectral ecology approach competitive with token methods
on timing.

\subsection{Joint Deconfounding}

The central unsolved problem is coverage deconfounding that works
for both intrinsic (curvature) and relational (cooperative stability)
invariants simultaneously.  v7's Sinkhorn regularization partially
deconfounds cooperative stability; v9's distribution comparison
partially deconfounds curvature.  But no single deconfounding strategy
addresses both, and combining separately-deconfounded components (as
v10 does) does not eliminate the joint confound.  A principled approach
might involve:

\begin{itemize}
\item Learning a coverage-invariant embedding space in which both
  curvature distributions and cooperative stability can be computed.
\item Using a permutation-based null model to calibrate each
  component against coverage-matched random pairs.
\item Constructing a joint likelihood model that explicitly
  includes coverage as a nuisance parameter.
\end{itemize}

\subsection{Beyond Module Abstraction}

The module abstraction layer---grouping k-mers into connected components
based on co-occurrence---may lose fine-grained structure.  Future work
could operate directly on weighted k-mer co-occurrence graphs, using
community detection to identify soft modules rather than hard
connected components.  This would preserve within-module heterogeneity
that the current binary assignment discards.


% ════════════════════════════════════════════════════════════════════
\section{Conclusion}
\label{sec:conclusion}

Spectral Ecology v10 is the culmination of a ten-version research
program that has systematically explored structural invariants of
genomic module ecosystems for reference-free relatedness detection.
The program began with the hypothesis that inherited genomic
organization can be detected through spectral analysis of module
co-occurrence graphs, and over ten iterations refined this hypothesis
through three eras of algorithmic development.

\begin{pullquote}
Ten versions.  Three eras.  One persistent confound.\\
The gap trajectory: $-26.87\% \to -0.61\% \to -1.66\%$.
\end{pullquote}

\textbf{What the program achieved:}

\begin{enumerate}
\item \textbf{Demonstrated positive separation.}  Eight of ten versions
  achieved positive separation, confirming that spectral invariants of
  module co-occurrence graphs contain genuine kinship signal.

\item \textbf{Approached the passing threshold.}  v7 achieved $-0.61\%$
  gap, within one percentage point of passing.  v10 achieved $-1.66\%$
  gap with better separation.  The gap trajectory shows consistent
  improvement from $-26.87\%$ through the gap-improving subsequence.

\item \textbf{Established complementary invariants.}  v9's intrinsic
  curvature and v7's cooperative stability detect different aspects of
  inherited structure, and their combination in v10 achieves a better
  separation--gap tradeoff than either alone.

\item \textbf{Identified the fundamental bottleneck.}  The coverage
  confound is not an algorithmic artifact---it is a structural property
  of module co-occurrence graphs that affects all spectral invariants.
  No combination of known invariants can eliminate it without explicit
  deconfounding.

\item \textbf{Produced correct rankings.}  v10 places all four
  parent--child pairs in the top 7 of 10 pairs, with three in the top 5.
  The ranking is informative even when the gap is negative.
\end{enumerate}

\textbf{What the program did not achieve:}

\begin{enumerate}
\item \textbf{Positive gap.}  No version of Spectral Ecology has
  achieved a positive weakest-pair gap.  Module Persistence v3
  ($+0.13\%$) and Coalition Transfer v3 ($+0.06\%$) pass using
  the same data and compute budget, operating at lower abstraction
  levels.

\item \textbf{Coverage deconfounding.}  Despite ten versions of
  progressively more sophisticated invariants, the coverage confound
  persists.  The spouse pair ranks in the top three in every version.

\item \textbf{Competitive performance.}  Rare-Run P90 achieves
  $+583\%$ separation and $+300\%$ gap---three orders of magnitude
  better than any spectral method.  The spectral approach is justified
  only if it provides insight that token methods cannot.
\end{enumerate}

\textbf{The honest assessment:}  Spectral Ecology is a research
program, not a production method.  It has demonstrated that structural
invariants of module co-occurrence graphs---eigenvalue distributions,
heat kernel signatures, curvature tensors, transport plans---contain
kinship signal that is detectable but not separable with current
techniques.  The program's value lies not in its results (which are
surpassed by simpler methods) but in its exploration of a fundamentally
different approach to genomic comparison: characterizing inherited
\emph{organization} rather than inherited \emph{content}.

Whether this approach can be made practical depends on solving the
coverage deconfounding problem, which v10 frames precisely but does
not solve.  The hybrid architecture provides the correct template---
combining intrinsic invariants with relational consistency measures,
weighted by discriminative power---but the components themselves must
be made coverage-invariant before the gap can turn positive.

The Spectral Ecology program continues.

\bibliographystyle{plain}
\bibliography{../references}

\end{document}
